\documentclass[../report_main.tex]{subfiles}
\begin{document}

\section{Background and Motivation}
With the rapid urbanization and the continuous increase in the number of vehicles on the road, traffic congestion has become a critical issue in modern cities worldwide. \gls{ITS} require robust and scalable simulation tools to evaluate traffic management strategies before they are implemented in the real world. A microscopic traffic simulation, which models the behavior and dynamics of individual vehicles rather than viewing traffic as a continuous flow, is essential for analyzing the complex interactions that lead to congestion, accidents, and inefficiencies. This project focuses on developing such a microscopic smart traffic simulation to model individual vehicle dynamics and their interactions within a highly customizable road network. Several existing projects and frameworks, such as Jraffic \cite{aminehabchi2023jraffic}, Java Traffic Simulation Systems \cite{zeeshier2023traffic}, and microscopic road traffic models like Movsim \cite{treiber2023movsim} and On-ramp simulations \cite{treiber2023onramp}, have laid the groundwork for this domain. Drawing inspiration from these, our project emphasizes an extensible, academic approach.

\section{Project Objectives}
The primary objective of this project is to construct a comprehensive Java-based traffic simulation application that adheres to \gls{OOP} principles, as taught in the university curriculum \cite{hustSlide}. The system is designed to simulate the movement of various vehicle types, including cars, buses, trucks, and emergency vehicles, by applying foundational physics equations for velocity and position updates. Furthermore, the simulation aims to implement highly realistic driver behaviors, encompassing normal, aggressive, and cautious driving profiles, which react dynamically to environmental stimuli such as leading vehicles and traffic lights. The project also targets the modeling of complex traffic networks with interconnected junctions, synchronized traffic lights, and multi-lane roads. A secondary but equally important objective is the development of a \gls{GUI} using JavaFX to visualize traffic flow in real-time and provide users with the capability to edit the map layout interactively. The application supports various modes, including a Map Editor mode for structural layout creation, a Simulation mode for observing traffic dynamics, and a Zoom mode for granular analysis.

\section{Report Structure}
The remainder of this report is organized systematically to detail the development and evaluation of the smart traffic simulation. Chapter 2 provides an in-depth exploration of the theoretical foundation, including vehicle kinematics and geometric collision detection. Chapter 3 focuses on the system architecture, detailing the \gls{OOP} design and the separation of core logic from the rendering engine. Finally, Chapter 4 concludes the report by summarizing the key achievements of the project and outlining potential avenues for future research and enhancement.

\end{document}